\begin{statement}{C.2.RAT.7}{Proved}{Every primitive rational double ordinary-triple cover has nonnegative Euler number}
Retain the torus-fibre, rational-normalization and specified
ruled-boundary morphism of C.2.RAT.1. Suppose $C$ is not
contracted, its image $H$ has generic ordinary-triple
transversal, and $C\to H^\nu$ has degree two. Suppose
the cubic normal tensor along $H$ has no common zero of
its coefficients. Then $e(S)\ge0$. This holds for every
splitting type of the normal bundle $N_{H/Y}$.
\end{statement}
\begin{proof}
Statement A.5.19 makes $H$ smooth: over a singular
rational ordinary-triple image every connected normalized
conductor cover has degree one. Thus $H\simeq\PP^1$.
Set $N=N_{H/Y}$, $W=\PP(N)$, and
$\xi=c_1(\OO_W(1))$, where $\PP(N)$ parameterizes lines.
The ambient Chern-class hypothesis gives
$\deg N=-2$ and $\OO_Y(S)|_H\simeq\OO_H$.
The cubic is an actual section of $\OO_W(3)$.
The assumed absence of a common coefficient zero makes
its root divisor finite of degree three over $H$:
\[
 R=R_1+R_2\sim3\xi,\qquad
 \xi^2=2,\qquad K_W=-2\xi,\qquad p_a(R)=4.
\tag{C.RD.11}
\]
Here $R_1$ is a section and hence is smooth rational;
$R_2$ has normalization $C\simeq\PP^1$, by the exact
integral-closure comparison (C.RD.2). Both are reduced
and they are different. The normalization sequence gives
\[
                  \sum_p\delta_p(R)=5.
\tag{C.RD.12}
\]
This equality also holds if the two components are
disjoint: the arithmetic-genus convention is
$p_a(R)=1-\chi(\OO_R)$ and the normalization has Euler
characteristic two.

Form the cyclic triple cover
$p_0:X_0\to W$ with algebra
$\OO_W\oplus\OO_W(-1)\oplus\OO_W(-2)$ and multiplication
given by the cubic. It is normal because its branch
divisor $R$ is reduced, and is connected because a
meridian around a smooth branch point permutes its
three sheets cyclically. Resolve its isolated
singularities and its reduced boundary to obtain
$X$, with $D=X\setminus U$, where
$U=\{v\in N:q_3(v)=1\}$ is the smooth normal-cone
level surface. The isomorphism over $W\setminus R$
is the exact radial-vector/dual-tautological-coordinate
comparison (C.D4.6).

The projective-bundle cohomology calculation gives
\[
 H^1(W,\OO_W)=H^1(W,\OO_W(-1))
             =H^1(W,\OO_W(-2))=0.
\tag{C.RD.13}
\]
Indeed both relative cohomologies of $\OO_W(-1)$
vanish, while $p_*\OO_W(-2)=0$ and
$R^1p_*\OO_W(-2)=\det N=\OO_{\PP^1}(-2)$.
Thus $H^1(X_0,\OO_{X_0})=0$ for every normal-bundle
splitting. Resolution Leray and C.2.RAT.6 give
$q(X)\le\sum p_g(X_0,x)\le1$, so $b_1(X)\le2$.

We spell out the boundary calculation with its number
of connected components retained. Write
\[
 \nu=\sum_p(r_p(R)-1),\qquad
 c=b_0(R),\qquad
 \lambda=\sum_x b_1(\operatorname{link}(X_0,x))\le2.
\]
Normalization of $R$ gives $e(R)=4-\nu$.
Since $U\to W\setminus R$ is etale of degree three,
$e(U)=3\nu$. The resolved boundary has the same
number $c$ of connected components. Replacing each
singular point of $R$ by the connected exceptional
divisor of its cubic suspension gives, by the
Mayer--Vietoris sequence with its individual branch
attachment points,
\[
                 b_1(D)=\nu-2+c+\lambda.
\tag{C.RD.14}
\]
The exceptional first Betti number is the link's
first Betti number: the exceptional intersection
matrix in its plumbed neighbourhood is invertible
over $\Q$. This is precisely the local comparison
proved with (C.D4.7). Additional boundary-resolution
blowups add rational trees and leave (C.RD.14)
unchanged.

Let $d_D$ be the number of irreducible components
of $D$, and let $V\subset H_2(X,\Q)$ be the span
of their classes. The exceptional classes are
independent. Pushing a relation to $W$ and pulling
back a relation from $R_1,R_2$ shows that
\[
 \varepsilon:=d_D-\dim V
      =2-\dim\operatorname{span}\{[R_1],[R_2]\}
      \in\{0,1\}.
\]
This proves the relation count rather than
assuming the two branch classes are independent.
The space $V$ contains the pullback of $R$,
whose square is $3R^2=54>0$. The Hodge index
theorem therefore makes its intersection pairing
nondegenerate. The divisor Gysin sequence gives
$b_1(U)=b_1(X)+\varepsilon\le3$; its degree-two
part identifies the image of $H_2(U)$ in
$H_2(X)$ with $V^\perp$. Consequently the rank
of the compact intersection pairing on $U$ is
$b_2(X)-d_D+\varepsilon$.
Now use $e(X)=2-2b_1(X)+b_2(X)$ and
$e(D)=c-b_1(D)+d_D$. Equations (C.RD.14) and
$e(U)=3\nu$ yield the exact identity
\[
\begin{split}
 \operatorname{rank}I_U
   &=e(U)-2+2b_1(X)+c-b_1(D)+\varepsilon\\
   &=2\nu+2b_1(X)-\lambda+\varepsilon
     \ge2\nu-2.
\end{split}
\tag{C.RD.15}
\]
In particular this formula includes disjoint root
components; no connected-boundary assumption has
been introduced.

If $\nu\ge2$, (C.RD.15) gives a nonzero compact
intersection pairing on a surface with $b_1(U)\le3$.
The compact-core transport (C.D4.2) embeds that
pairing into an actual nearby torus. Its induced
first-homology image has dimension at most three,
so its second-homology image has zero wedge pairing,
as established immediately after (C.D4.2). This
contradiction proves $\nu\le1$.

Finally suppose a value of $H$ receives vertical
boundary attachments at both its $C$-preimages.
If their tangent-root points in $R_2$ were distinct,
at least one would be a simple root of the cubic,
and (C.RD.3)--(C.RD.4) would give pure conductor
there, a contradiction. They therefore map to a
single multibranch point of $R_2$. Distinct values
give distinct such points, each contributing at
least one to $\nu$. Hence there is at most one
value with two attachments. With
$K\ge r-N_0-r_{\rm dom}$ surviving original
vertical boundaries, at least $K-1$ distinct
normalization-incidence edges remain at $H$.
There is only one remaining generic triple branch,
so $r_{\rm dom}\le1$. The graph comparison
(C.RA.24) and exact Euler decomposition now give
\[
 e(S)\ge2-r+N_0+(K-1)
                         \ge1-r_{\rm dom}\ge0.
\tag{C.RD.16}
\]
This completes the proof.
\end{proof}


% C.2.RAT.8 from part_C2_rational_triple.tex

